\documentclass[11pt]{article}
\usepackage[margin=0.9in]{geometry}
\usepackage{amsmath,amssymb,amsthm,booktabs}
\usepackage[T1]{fontenc}
\usepackage{lmodern,microtype}
\usepackage[colorlinks=true,linkcolor=blue,urlcolor=blue]{hyperref}
\newtheorem{theorem}{Theorem}[section]
\newtheorem{lemma}[theorem]{Lemma}
\newtheorem{proposition}[theorem]{Proposition}
\newtheorem{corollary}[theorem]{Corollary}
\theoremstyle{remark}
\newtheorem{remark}[theorem]{Remark}
\newcommand{\R}{\mathbb R}
\newcommand{\C}{\mathbb C}
\newcommand{\norm}[1]{\left\|#1\right\|}
\newcommand{\diag}{\operatorname{diag}}
\newcommand{\ran}{\operatorname{ran}}
\newcommand{\HH}{\mathcal H}
\title{Toward a four-dimensional Yang--Mills construction\\
\large Notebook 12: fine regularity on the flat current domain}
\author{Research working document}
\date{16 October 2026}
\begin{document}
\maketitle
\begin{abstract}
We strengthen the flat fine inverse from physical-cell Hilbert bounds
to pointwise regularity for arbitrary bounded fine forcing. A discrete
heat-kernel argument and the canceled finite-rank formula give uniform
$C^{1,\alpha}$ estimates for every $\alpha<1$, exponential cell locality,
and a componentwise Laplacian bound at zero current. An explicit
counterexample rules out the endpoint and individual second-difference
bounds. Scalar gauge identities then place the full current
perturbation in the actual homogeneous source norms $N$ and $N_\beta$,
without differentiating the current. A smaller uniform complex domain
supports the full flat inverse and extension on bounded forcing.
Right-divergence data, nonflat backgrounds, and retained histories
require additional estimates stated separately below.
\end{abstract}

\noindent\textbf{Status.}
Research proofs checked by separate internal derivations and audits.
No external certification or novelty claim is made. The baseline is
valid for every odd $M$. The nonlinear constraint-current term uses
$M=L^j$ with fixed odd $L\ge3$, as proved in Notebook 11. All norms
use physical fine derivatives and volume, at final block scale one.

\section{Previously established algebraic inputs}

We use the exact homogeneous inverse and scalar source definitions of
Notebook 09, and the corrected quadratic-composition bound of Notebook 11.
Let $D$ be the fine scalar gradient, $\mathcal L=D^*D$, $Q_0$ the
normalized scalar block average, and $Q=Q_M^{\rm st}$ the normalized
straight one-form average. The scalar gauge is
\[
 G'=(\mathcal L+Q_0^*a'Q_0)^{-1},\qquad
 R=I-G'Q_0^*(Q_0G'^2Q_0^*)^{-1}Q_0G',\qquad Z=G'RG'.
\]
The distinct positive penalty $a$ enters the fine one-form inverse.
The scalar penalty $a'$ does not replace it. The identities
\[
 Q_0Z=0,\quad \mathcal L G'R=R,\quad
 Z\mathcal L=G'R,\quad \mathcal L Z\mathcal L=R
\]
were proved from these definitions. They will also justify the
fixed-refinement operator products on bounded forcing.

The canceled Bloch formulas $G_0=F+\beta E_a^{-1}\beta^-$,
including every central entry, are those proved in Notebook 09.
We reuse their positive common strip and uniform coefficient bounds;
no new pointwise assertion is assumed from their Hilbert estimates.
For the scalar notation below, with $p_r=z+2\pi r$, set
\[
 a_r=\prod_\nu M^{-1}\sum_{s\in I_M}e^{ip_{r,\nu}s/M},\qquad
 \omega_{r,\mu}=a_rM^{-1}\sum_{t=0}^{M-1}e^{ip_{r,\mu}t/M},
\]
\[
 R_{s,2}=\sum_{r\ne0}\frac{a_ra_r^-}{P_r^2},\qquad
 D_{s,2}=a_0a_0^-+P_0^2R_{s,2}.
\]
The symbols $P_r$ and the negative-argument convention are recalled
in the next section. These are scalar denominators; the subscript
$s,2$ is a name and does not designate a free matrix index.
The homogeneous source value of the one-form penalty is $a=1$;
a fixed compact positive interval is allowed throughout.
% Pointwise flat inverse; no standalone preamble.
% Requires amsmath/amssymb/amsthm, theorem environment, and \norm.
% Every label introduced by this fragment has the unique prefix pt-.
\section{Pointwise regularity for arbitrary fine forcing}
\label{pt-section}

The full flat inverse also obeys pointwise estimates. These estimates
concern arbitrary bounded fine forcing, without assuming that it is a
coarse lift. Retain fine spacing $\eta=M^{-1}$, odd $M$, unit coarse
cells, the physical fine measure $\eta^d$, and
\begin{equation}\label{pt-operator}
 \mathcal A_a=d^*d+DRD^*+Q^*aQ,\qquad G_0=\mathcal A_a^{-1},
 \qquad \mathcal L=D^*D.
\end{equation}
Here $d$ is the fine curl, $D$ the scalar gradient, $Q=Q_M^{\rm st}$,
and $R$ is the matched homogeneous flat scalar source projection.
The penalty lies in a fixed interval $0<a_-\le a\le a_+<\infty$;
the normalized source value $a=1$ is included. The discussion applies
to the homogeneous infinite lattice and compatible periodic regulators.

For a fixed coarse-cell patch $U$, define the discrete pointwise norm
\begin{equation}\label{pt-holder-norm}
 \norm{u}_{1,\alpha;U}=\norm{u}_{\infty;U}
  +\sum_\mu\norm{D_\mu^+u}_{\infty;U}
  +\sum_\mu\sup_{\substack{x,y\in U\\x\ne y}}
       \frac{|D_\mu^+u(x)-D_\mu^+u(y)|}{|x-y|_1^\alpha}.
\end{equation}
The points are actual fine-grid points and a fixed collar supplies the
differences. Enlarging the patch by finitely many neighboring coarse
cells changes only constants. Backward differences obey the same bounds.

\begin{theorem}\label{pt-theorem}
For each $0<\alpha<1$ there are $C_\alpha,\rho>0$, independent of
odd $M$ and compatible periodic volume, such that the block from
forcing in coarse cell $m$ to a fixed collar of coarse cell $n$ satisfies
\begin{align}
 \norm{G_0(n,m)}_{\ell^\infty\to C^{1,\alpha}}
       &\le C_\alpha e^{-\rho|n-m|_1},
       \label{pt-pointwise-bound}\\
 \norm{(\mathcal L G_0)(n,m)}_{\ell^\infty\to\ell^\infty}
       &\le C e^{-\rho|n-m|_1}.
       \label{pt-laplacian-bound}
\end{align}
The first target norm is \eqref{pt-holder-norm}; $\mathcal L$ in
the second estimate acts componentwise on the vector output.
On tori use periodic coarse-cell distance. Weighted row and column
block sums are uniformly bounded at every smaller positive exponent.
Consequently bounded fine forcing produces uniformly bounded fields,
first differences, locally $\alpha$-H\"older first differences, and
componentwise Laplacians.
\end{theorem}

The constants need not stay bounded as $\alpha\uparrow1$. In fact a
uniform Lipschitz estimate for the first differences, and a uniform
bound for each individual second difference, are false in dimension
$d\ge2$; an example is given below. No source weak-norm identification
is inferred merely from the theorem's terminology.

\subsection{Reduction to a scalar high-mode inverse}

Use the aliases $p_r=z+2\pi r$, $r\in I_M^d$, with
$I_M=\{-(M-1)/2,\ldots,(M-1)/2\}$, and write
\begin{equation}\label{pt-alias-data}
 u_{r,\mu}^{\pm}=M(e^{\pm ip_{r,\mu}/M}-1),\qquad
 P_r=\sum_\mu u_{r,\mu}^+u_{r,\mu}^-,\qquad
 \mathcal F_M(z)_{sr}=e^{i(z+2\pi r)\cdot s/M}.
\end{equation}
The common holomorphic chart from the full inverse theorem contains
the real fundamental cube in the interior of $|\operatorname{Re}z_\mu|<K<2\pi$.
For real $z$, $\mathcal F_M$ is unitary from alias counting $\ell^2$
to physical-cell $L^2$. On the complex chart it and its inverse have
uniformly bounded $L^2$ norms. In particular a bounded fine-cell input
has alias coefficients bounded in $\ell^2$, and individually, by
$C\norm{f}_\infty$.

Recall the canceled inverse $G_0=F+\beta E_a^{-1}\beta^-$.
Separate the scalar orientation-diagonal multiplier
\begin{equation}\label{pt-high-scalar}
 T_r=P_r^{-1}\quad(r\ne0),\qquad T_0=0.
\end{equation}
The difference $F-T$ is a fixed sum of rank-one operators per vector
orientation: its central block, two mixed blocks, and one high/high
term. Its high vectors are $\omega_{r,\mu}^-/P_r$ and the analytic
transpose partners. The remaining scalar coefficients are uniformly
bounded and holomorphic. Thus the only full-rank pointwise problem is
\eqref{pt-high-scalar}. Minus signs on multipliers throughout this
argument denote evaluation at negative full momentum, with the
analytic transpose; they do not conjugate a complex variable.

\subsection{Weighted heat estimates with complex Bloch twist}

Let $p_s(n)$ be the continuous-time nearest-neighbor heat kernel on
$\mathbb Z$ with generator $f(n+1)+f(n-1)-2f(n)$. Its Poisson-process
representation gives
\begin{equation}\label{pt-heat-moment}
 \sum_n e^{tn}p_s(n)=e^{2s(\cosh t-1)},\qquad
 \sum_n p_s(n)=1.
\end{equation}
It is symmetric and decreasing for $n\ge0$. For completeness, set
$q_n=p_s(n)-p_s(n+1)$ on $n\ge0$. Then
\[
 q_0'=q_1-3q_0,\qquad
 q_n'=q_{n-1}+q_{n+1}-2q_n\quad(n\ge1),
\]
with nonnegative initial data. This bounded cooperative system
preserves nonnegativity, as follows by shifting its generator to a
matrix with nonnegative entries and expanding the exponential.
Also Fourier inversion gives
\begin{equation}\label{pt-heat-origin}
 p_s(0)=\frac1{2\pi}\int_{-\pi}^{\pi}
               e^{-4s\sin^2(t/2)}\,dt
       \le C(1+s)^{-1/2}.
\end{equation}
For $0\le r\le r_0$, summation by parts therefore yields the exact
weighted variation identity
\begin{align}
 \sum_n e^{r|n|}|p_s(n+1)-p_s(n)|
  &=(1+e^r)\left[p_s(0)
       +(1-e^{-r})\sum_{n\ge1}e^{rn}p_s(n)\right]\notag\\
  &\le C\left[(1+s)^{-1/2}
                   +r e^{2s(\cosh r-1)}\right].
       \label{pt-weighted-variation}
\end{align}
Take $s=M^2t$ and $r=\delta/M$, where $\delta$ is a fixed small
constant. Multiplying the difference by $M$ gives
\begin{equation}\label{pt-first-heat}
 \sum_n e^{\delta|n|/M}
       |D_\eta p_{M^2t}(n)|
       \le C(t^{-1/2}+\delta)e^{C\delta^2t}.
\end{equation}
For $t<\eta^2$ the trivial $M$ bound is smaller than $t^{-1/2}$.
The weight is submultiplicative. Apply
\eqref{pt-first-heat} twice to
$p_s=p_{s/2}*p_{s/2}$ to obtain the second-difference bound.
Tensor products then give, for any specified composition of $k$ fine
forward/backward differences in $d$ dimensions,
\begin{equation}\label{pt-heat-two-derivatives}
 \sum_{n\in\mathbb Z^d}e^{\delta|n|_1/M}
       |D_\eta^k p_{M^2t}(n)|
       \le Ct^{-k/2},\qquad 0<t\le1,\quad k=0,1,2.
\end{equation}
The displayed $p$ is a probability kernel. If it is instead written
as a density against physical fine volume, the density factor $M^d$
cancels the summation factor $M^{-d}$.

The reconstructed Bloch heat operator has the exact image-sum kernel
\begin{equation}\label{pt-twisted-heat}
 H_z(t;s,s')=\sum_{k\in\mathbb Z^d}
             e^{-iz\cdot k}p_{M^2t}(s-s'+kM).
\end{equation}
This convention uses ordinary summation over the fine-cell input
index $s'$. The formula follows by finite Fourier inversion and
continues to complex $z$ by the exponential moment bound.
For $s$ in a fixed output collar and $s'$ in the input cell,
\[
 |e^{-iz\cdot k}|
       \le C e^{\delta|s-s'+kM|_1/M}.
\]
Thus \eqref{pt-heat-two-derivatives}, including the corresponding
output differences of the image sum, proves
\begin{equation}\label{pt-twisted-derivatives}
 \norm{D_\eta^kH_z(t)}_{\infty\to\infty}
       \le Ct^{-k/2},\qquad 0<t\le1,\quad k=0,1,2.
\end{equation}
The difference between first-gradient rows at fine points $x,y$ is
bounded both by $Ct^{-1/2}$ and, by telescoping, by
$C|x-y|_1t^{-1}$. Interpolating these two elementary bounds gives
\begin{equation}\label{pt-heat-holder}
 \norm{H_z(t)}_{\infty\to C^{1,\alpha}}
       \le C_\alpha t^{-(1+\alpha)/2},\qquad 0<t\le1.
\end{equation}
Subtracting the central heat mode changes this only by a bounded
term, since $P_0,u_0$ and its plane-wave profile are uniformly bounded
on the chart. For $t\ge1$, the high-alias bound
$\operatorname{Re}P_r\ge c|r|^2$ and the individual input-coefficient
bound instead give
\begin{align}
 &\norm{H_z(t)-e^{-tP_0}\Pi_0(z)}_{\infty\to C^{1,\alpha}}
       \notag\\
 &\qquad\le C\sum_{r\ne0}(1+|r|)^{1+\alpha}e^{-ct|r|^2}
       \le C_\alpha e^{-ct/2}.
       \label{pt-heat-large-time}
\end{align}
The high-mode inverse is therefore the convergent operator integral
\begin{equation}\label{pt-scalar-integral}
 T(z)=\int_0^\infty
             \bigl[H_z(t)-e^{-tP_0}\Pi_0(z)\bigr],dt.
\end{equation}
The small-time integral is finite because
$\int_0^1t^{-(1+\alpha)/2}\,dt=2/(1-\alpha)$.
This proves a uniform holomorphic
$\ell^\infty\to C^{1,\alpha}$ bound for $T$ for every $\alpha<1$.

\subsection{Pointwise bounds for the rank corrections}

The centered scalar and straight-vector averaging multipliers satisfy
\begin{equation}\label{pt-alias-weights}
 |a_r^\pm|\le Cw(r),\qquad
 |\omega_{r,\mu}^\pm|\le\frac{Cw(r)}{1+|r_\mu|},\qquad
 w(r)=\prod_\nu(1+|r_\nu|)^{-1}.
\end{equation}
Indeed, for a nonzero component, the numerator of
$\sin((z_\nu+2\pi r_\nu)/2)/
 [M\sin((z_\nu+2\pi r_\nu)/(2M))]$
is uniformly bounded and its denominator has modulus at least
$c|r_\nu|$ on the enlarged chart. The zero component is bounded
directly. The vector multiplier has one extra longitudinal factor.
This includes aliases at the edge of the fine zone.

On a dyadic maximum-norm shell,
\begin{equation}\label{pt-shell-sum}
 \sum_{2^j\le|r|_\infty<2^{j+1}}w(r)
                      \le C_d(j+1)^{d-1}.
\end{equation}
Choose a coordinate at the maximal scale. Its reciprocal weight sums
to $O(1)$ over that shell, and each other coordinate contributes
$O(j+1)$. Summing over the possible maximal coordinates proves
\eqref{pt-shell-sum}. Consequently
\begin{equation}\label{pt-rank-holder-sum}
 \sum_{r\ne0}(1+|r|)^{1+\alpha}
             \left|\frac{\omega_{r,\mu}^-}{P_r}\right|
   \le C\sum_{r\ne0}w(r)(1+|r|)^{\alpha-1}
   \le C_\alpha.
\end{equation}
Absolute Fourier summation now bounds the reconstructed high rank
vectors in $C^{1,\alpha}$. Their central partners have the same bound.
The Woodbury vector has
\begin{equation}\label{pt-beta-bound}
 |\beta_{r,\mu}|
  \le C\left(\frac{|a_r^-|}{(1+|r|)^3}
                  +\frac{|\omega_{r,\mu}^-|}{(1+|r|)^2}\right),
       \qquad r\ne0,
\end{equation}
so its output profile satisfies the same weighted sum. Its canceled
central component is bounded directly.

The input rows of every rank term have uniformly bounded alias
$\ell^2$ norm. By the physical-cell Fourier normalization,
\[
 |\text{input row}(f)|
       \le C\norm{f}_{L^2(\text{physical cell})}
       \le C\norm{f}_\infty.
\]
The cell has physical volume one; there is no factor $M^{d/2}$.
Combining these row estimates with
\eqref{pt-rank-holder-sum}--\eqref{pt-beta-bound} controls all rank
corrections. The full reconstructed $G_0$ therefore has the required
uniform holomorphic pointwise bound on the common chart.

\subsection{The componentwise Laplacian}

For the scalar term, $\mathcal L T=I-\Pi_0$ is bounded on
$\ell^\infty$. Applying $\mathcal L$ to a high rank vector
$\omega_{r,\mu}^-/P_r$ gives $\omega_{r,\mu}^-$. Its full
reconstruction is a column of the physical adjoint average, minus
its bounded central mode. Explicitly, for $x=n+s/M$, $M=2h+1$,
\begin{equation}\label{pt-adjoint-average}
 (Q^*B)_\mu(x)=b_s B_\mu(n)+(1-b_s)B_\mu(n-e_\mu),
 \qquad b_s=\frac{s_\mu+h+1}{M}.
\end{equation}
These weights are nonnegative and sum to one. The Bloch profile is
$b_s+(1-b_s)e^{-iz_\mu}$, uniformly bounded on the complex chart
and any fixed collar. An absolute alias sum of $\omega$ is unnecessary
and would introduce a spurious logarithmic loss.

For the Woodbury vector, its high-alias formula gives
\begin{equation}\label{pt-beta-laplacian}
 (\mathcal L\beta)_{r,\mu}
  =\frac{P_0u_{r,\mu}^+a_r^-}{P_rD_{s,2}}
           -\omega_{r,\mu}^-k_\mu d_\mu^+.
\end{equation}
The first term is $O(w(r)/(1+|r|))$, absolutely summable by
\eqref{pt-shell-sum}. The second is a bounded coefficient times
the profile \eqref{pt-adjoint-average}, with its central subtraction.
The canceled central Laplacian component is bounded. The same input
row estimates as before now prove a uniform holomorphic
$\ell^\infty\to\ell^\infty$ bound for $\mathcal L G_0$.
This argument supplies neither H\"older continuity of its Laplacian
nor bounds on all separate second differences. In particular the
$Q^*$ profile can have transverse block jumps.

\subsection{Gluing, exponential decay, and bounded forcing}

The whole reconstructed inverse and its fine derivatives are periodic
in coarse momentum. The full real symbol is invariant under the
alias relabeling induced by $z\mapsto z+2\pi n$; the within-cell
integer phase cancels wrapping of the alias index. The enlarged
holomorphic charts agree on their real overlaps and hence glue.
Individual separated central/high terms need not be periodic; they
were used only to bound the whole symbol on each chart.

Output collars consist of finitely many neighboring-cell copies,
with bounded holomorphic factors $e^{iz\cdot n}$. Thus the common
strip bounds hold as operator bounds into the discrete norm
\eqref{pt-holder-norm}. For each fixed $M$ these are finite-dimensional
Banach spaces, so entrywise holomorphy and the proved uniform norms
give Banach-valued contour shifting. This proves
\eqref{pt-pointwise-bound}--\eqref{pt-laplacian-bound} without a
growing-dimension factor.

In infinite volume the resulting absolutely summable blocks define
the inverse on bounded forcing and agree with the Hilbert inverse
on finitely supported fine forcing. The inverse identity extends
through absolute convergence of the exponentially summable operator
products at each fixed $M$. The source projection is not local, but
$R=\mathcal L Z\mathcal L$ has exponentially summable cell $L^2$
blocks. Finite-cell norm equivalence makes their pointwise blocks
summable for each fixed $M$. That last identity argument need not
be uniform in $M$: the pointwise inverse bound was proved uniformly
by the heat and rank estimates above.

Finite compatible tori sample the same symbol, giving exact
periodizations and the same bounds in periodic distance, including
short periods and multiplicities. Summing the cell kernels bounds
the field and gradient globally and the gradient H\"older seminorm
on every fixed collar. Nearby points lie in one enlarged collar;
farther points use twice the gradient supremum. This also gives the
global discrete H\"older bound, with periodic physical distance on
tori.

\subsection{A sharp obstruction to the Lipschitz endpoint}

For $d\ge2$, take the compatible torus consisting of one coarse
cell. At its sole coarse momentum, zero, the full inverse is
$a^{-1}$ on the central constant mode and $P_r^{-1}$ on every high
alias. Thus on fine mean-zero forcing it is exactly the scalar
periodic inverse Laplacian, independently in each vector component.

Choose a smooth cutoff $\chi$ supported in a small ball inside the
unit continuum torus and equal to one near its center. Set
\begin{equation}\label{pt-endpoint-example}
 u_\varepsilon(x)=\chi(x)x_1x_2\,
                       \frac12\log(|x|^2+\varepsilon^2).
\end{equation}
For $P=x_1x_2$ and $f(r)=\frac12\log(r^2+\varepsilon^2)$,
harmonicity and homogeneity of $P$ give
\begin{equation}\label{pt-endpoint-laplacian}
 \Delta(Pf)=P\left(f''+\frac{d+3}{r}f'\right)
   =P\frac{(d+2)r^2+(d+4)\varepsilon^2}
                  {(r^2+\varepsilon^2)^2}.
\end{equation}
Since $|P|\le r^2/2$, this is bounded independently of $\varepsilon$.
The cutoff-region derivatives are also uniformly bounded because
that region stays away from zero. However
\begin{equation}\label{pt-endpoint-divergence}
 \partial_1\partial_2u_\varepsilon(0)=\log\varepsilon.
\end{equation}

For fixed $\varepsilon>0$, sample the smooth periodic function on
increasingly fine odd grids and subtract its exact discrete mean.
Define $f_{M,\varepsilon}=\mathcal L_\eta u_{M,\varepsilon}$.
Its mean is zero, and Taylor convergence gives
$\norm{f_{M,\varepsilon}}_\infty\le C+1$ for sufficiently large
$M$, depending on $\varepsilon$. The full inverse returns exactly
$u_{M,\varepsilon}$. Its mixed second fine difference at the origin
converges to \eqref{pt-endpoint-divergence}. Choosing small
$\varepsilon$ and then large $M$ proves that bounded forcing can
produce arbitrarily large mixed second differences. A uniform
Lipschitz bound on first differences would bound these differences
and is therefore impossible. This does not contradict
\eqref{pt-laplacian-bound}, which controls their Laplacian sum rather
than each second difference separately.

\subsection{Source scope}

The theorem supplies arbitrary-forcing pointwise estimates for the
matched homogeneous flat inverse, including its normalized $a=1$
case. It does not by itself identify every numerical factor or term
in the source's weak norms, establish nonflat estimates, handle
general retained-history domains, or control the complete interacting
current carrier. Those applications require exact norm and coordinate
checks. The false endpoint above must not enter such a match through
an implicit assertion that $C^{1,\alpha}$ controls all second derivatives.
% Pointwise full-current neighborhood; no standalone preamble.
% Requires amsmath/amssymb/amsthm, theorem environment, and \norm.
% Uses the preceding pointwise baseline theorem. Labels have prefix pc-.
\section{The full flat current in the pointwise weak norms}
\label{pc-section}

We now transfer the pointwise estimates to arbitrary bounded complex
fine currents. The background remains exactly flat. For the corrected
nonlinear constraint term, refinement is $M=L^j$ with fixed odd $L$,
as in the independently audited fixed-step quadratic-composition
theorem. The result does not use a pointwise bound for the projection
$\Pi$ on arbitrary bounded fields and does not differentiate the current.

\subsection{The source norms and spatial weights}

At homogeneous $\ell=1$ and background $U=1$, deposited Raw I,
lines 12740--12752, defines the equivalent norms
\begin{equation}\label{pc-source-norms}
 N(A)=\max\{\norm{A}_\infty,\norm{\nabla A}_\infty\},\qquad
 N_\beta(A)=\max\{N(A),[\nabla A]_\beta\},
 \quad 0\le\beta\le\beta_0<1.
\end{equation}
Here $\nabla$ collects the first fine differences of a vector field;
the H\"older member can be taken on fixed coarse-cell collars.
The pointwise estimates of the preceding section also imply its
global discrete version. There is no Laplacian in $N_\beta$.
The source explicitly excludes the Laplacian member for the current
inverse at Raw I, lines 13759--13765.

Let $\varpi(n)>0$ be a cell weight satisfying
\begin{equation}\label{pc-weight-ratio}
 \frac{\varpi(n)}{\varpi(m)}
       \le e^{\kappa\operatorname{dist}(n,m)}.
\end{equation}
The distance is ordinary or periodic cell distance. Define
$\norm{f}_{\infty,\varpi}$, $N_\varpi$, and $N_{\beta,\varpi}$
by the corresponding weighted supremum over cells and fixed collars.
Neighboring weights have uniformly bounded ratios. In particular
$\varpi(n)=e^{\kappa\operatorname{dist}(n,m_0)}$ is allowed.
These norms are applied to complete fields. We never sharply cut a
differentiable field into cells before applying $D^*$, which would
introduce artificial fine-scale boundary derivatives.

Retain
\begin{equation}\label{pc-scalar-data}
 \mathcal L=D^*D,\qquad Z=G'RG',\qquad
 \Lambda=Z\mathcal L D^*=G'RD^*,\qquad \Pi=I-D\Lambda.
\end{equation}
Stars denote physical adjoints on the real slice and their fixed
analytic bilinear continuations for complex parameters. The
penalized baseline is
$G_0=(d^*d+DRD^*+Q^*aQ)^{-1}$, with matched unit-scale value $a=1$.

\subsection{The scalar smoothing needed for the current}

For a common small $\kappa>0$, the scalar inverse supplies constants
independent of $M$ and compatible volume such that
\begin{align}
 \norm{Z\mathcal L f}_{\infty,\varpi}
       &\le c_0\norm{f}_{\infty,\varpi},&
 \norm{DZ\mathcal L f}_{\infty,\varpi}
       &\le c_1\norm{f}_{\infty,\varpi},
       \label{pc-scalar-bounded}\\
 N_{\beta,\varpi}(DZf)
       &\le z_1\norm{f}_{\infty,\varpi},&
 N_{\beta,\varpi}(DZD^*f)
       &\le z_2\norm{f}_{\infty,\varpi}.
       \label{pc-scalar-holder}
\end{align}
Constants can be chosen at the fixed ceiling $\beta_0$, uniformly
for $0\le\beta\le\beta_0$.

Here is the pointwise proof of these additional scalar facts. The
canceled alias matrix of $Z$ is the high diagonal $P_r^{-2}$ plus
a fixed number of rank terms whose high vectors are $a_r^-/P_r^2$
and their transpose partners. Its central entries have no massless
pole. The high diagonal has the heat representation
\begin{equation}\label{pc-biharmonic-heat}
 \int_0^\infty t\,[H_z(t)-e^{-tP_0}\Pi_0(z)]\,dt.
\end{equation}
The weighted heat derivative bounds extend to every fixed integer
order $k$: split the semigroup into $k$ equal-time factors, put one
fine difference on each, and use weighted $L^1$ convolution.
Thus the small-time bound is $C_kt^{-k/2}$. An additional
$\beta$-H\"older output difference costs $t^{-\beta/2}$ by the
same interpolation used for the baseline inverse. Input derivatives
are derivatives of the kernel, or discrete summation by parts, and
do not require derivatives of the forcing. The relevant integral is
\[
 \int_0^1 t^{1-(m+\beta)/2}\,dt<\infty
                  \quad\text{if }m+\beta<4.
\]
For $Z\mathcal L,DZ\mathcal L$ the derivative orders are two and
three. For the gradient H\"older members of $DZ,DZD^*$ they are
two and three, respectively, before the extra $\beta$. Large times
use the high-alias gap.

For the rank terms, an input row with the required derivatives has
uniform alias $\ell^2$ norm. Output vectors obey
$|a_r/P_r^2|\le Cw(r)/(1+|r|)^4$, where
$w(r)=\prod_\nu(1+|r_\nu|)^{-1}$ and its dyadic shell sum is
$O((j+1)^{d-1})$. Every output derivative/H\"older order strictly
below four is therefore absolutely summable. Physical-cell $L^2$
input rows act on bounded forcing without a dimension factor.
The same periodic holomorphic gluing as for $G_0$ gives exponential
cell kernels and hence \eqref{pc-scalar-bounded}--\eqref{pc-scalar-holder}.
Finite periodic sampling preserves the bounds.

It follows, using the actual derivative of the complete field, that
\begin{equation}\label{pc-lambda-on-N}
 \norm{\Lambda A}_{\infty,\varpi}\le l_0N_\varpi(A),\qquad
 \norm{D\Lambda A}_{\infty,\varpi}\le l_1N_\varpi(A),\qquad
 \norm{\Pi A}_{\infty,\varpi}\le(1+l_1)N_\varpi(A).
\end{equation}
No assertion that $D\Lambda$ or $\Pi$ is bounded on unrestricted
$\ell^\infty$ input is used.

\subsection{Both current terms map $N$ to $N_\beta$}

The matched local auxiliary factorization is
\begin{equation}\label{pc-aux-factorization}
 \Theta'_J=\mathfrak j_J^*\Lambda+\Lambda^*\mathfrak j_J
       +\Lambda^*V_JD\Lambda+(D\Lambda)^*V_J^*\Lambda.
\end{equation}
The finite-incidence coefficient operators $\mathfrak j_J,V_J$
and their analytic transposes have weighted pointwise norms at most
$c_j\norm{J}_\infty,c_v\norm{J}_\infty$. They contain no derivative
of $J$. The exact scalar identity $G_0\Lambda^*=DZ$ gives
\begin{equation}\label{pc-outer-aux-smoothing}
 G_0\Theta'_J
  =G_0\mathfrak j_J^*\Lambda+DZ\mathfrak j_J
       +DZV_JD\Lambda+DZD^*V_J^*\Lambda.
\end{equation}
Let $g_\beta$ bound $G_0$ from weighted bounded forcing to
$N_{\beta,\varpi}$. Then
\begin{equation}\label{pc-aux-N-bound}
 N_{\beta,\varpi}(G_0\Theta'_JA)
        \le C_{\rm aux}\norm{J}_\infty N_\varpi(A),
\end{equation}
where one possible constant is
\[
 C_{\rm aux}=g_\beta c_jl_0+z_1c_j+z_1c_vl_1+z_2c_vl_0.
\]
This is the homogeneous normalized source (J1), including weights.
Its input requires $N$ only, not H\"older regularity of $A$ or $J$.

For the full constraint term, write
\begin{equation}\label{pc-constraint-weight}
 \Theta_J=T_{v_J},\qquad v_J=E_J^T J,
 \qquad
 \langle A,T_vB\rangle_\eta
      =2\langle C_j^{(2)}(A,B),v\rangle_1,
\end{equation}
where $E_J$ is the auxiliary normalized extension and $C_j^{(2)}$
is the quadratic jet of the actual corrected fixed-$L$ composition.
The previously proved auxiliary cell-Schur estimate gives, analytically,
$\norm{v_J}_\infty\le C_v\norm{J}_\infty$ throughout its current ball.

The direct pointwise estimate for $T_v$ uses the same physical
factor as its Hilbert proof. At a step of fine spacing $h$, the
quadratic coefficient is $(h/L)C_L$. Its finite coefficient tensor,
paired with coarse measure and converted to the fine physical
adjoint, has an absolute row sum at most $c_Lh\norm{v}_\infty$.
The factor $L^d$ in this conversion is fixed. The quadratic chain
rule expresses the full operator as
\begin{equation}\label{pc-pointwise-chain}
 T_v=\sum_{t=0}^{j-1}Q_{0:t}^*T_{h_t,v_t}Q_{0:t},\qquad
 h_t=L^{t-j},\qquad v_t=Q_{t+1:j}^*v.
\end{equation}
Both flat $Q$ and its physical adjoint $Q^*$ have nonnegative row
coefficients summing to one. They are $\ell^\infty$ contractions,
including the covering-space multiplicities on short tori. Thus
\[
 \norm{T_v}_{\infty\to\infty}
       \le\frac{c_L}{L-1}\norm{v}_\infty.
\]
The composition has fixed unit-cell range $R_*$, since its reading
radii sum over the same geometric series. Consequently
\begin{equation}\label{pc-T-pointwise}
 \norm{T_vf}_{\infty,\varpi}
       \le\tau_\infty\norm{v}_\infty
                         \norm{f}_{\infty,\varpi}.
\end{equation}
The constant includes $e^{\kappa R_*}$. This is a coefficient-row
argument, not a conversion of a final-cell Hilbert norm into a
pointwise norm.

The identity
\begin{equation}\label{pc-Pi-outer}
 G_0\Pi^*=G_0-DZD^*
\end{equation}
now places every smoothing operator outside the bounded current
coefficients. Combining \eqref{pc-lambda-on-N},
\eqref{pc-T-pointwise}, and \eqref{pc-Pi-outer} gives
\begin{align}
 N_{\beta,\varpi}(G_0\Pi^*\Theta_J\Pi A)
       &\le C_{\rm con}\norm{J}_\infty N_\varpi(A),
       \label{pc-constraint-N-bound}\\
 C_{\rm con}&=(g_\beta+z_2)\tau_\infty C_v(1+l_1).
       \notag
\end{align}
In particular no derivative of $J$ has appeared.

There is also a constrained baseline estimate. Put
\begin{equation}\label{pc-constrained-baseline}
 \Gamma=G_0-H_{\rm src}q_{\rm aux}^{-1}H_{\rm src}^*,
 \qquad q_{\rm aux}=q_M+aI.
\end{equation}
The extension $H_{\rm src}$ is pointwise local into $N_\beta$;
its transpose maps bounded fine fields to bounded coarse fields
by the physical-cell Schur bound; and
$q_{\rm aux}^{-1}=QG_0Q^*$ is a local bounded coarse operator.
Hence $\Gamma$ maps weighted bounded forcing to $N_\beta$ and
\begin{equation}\label{pc-constrained-current}
 N_{\beta,\varpi}(\Gamma\Theta_JA)
      \le C\norm{J}_\infty\norm{A}_{\infty,\varpi}.
\end{equation}

\subsection{A common pointwise current neighborhood}

\begin{theorem}\label{pc-theorem}
Let $\varepsilon_H$ be the established full-current Hilbert/Schur
radius, which is contained in the auxiliary radius. Define
\begin{equation}\label{pc-radius}
 C_*=C_{\rm aux}+C_{\rm con},\qquad
 \varepsilon_{\rm pt}
       =\min\{\varepsilon_H,(2C_*)^{-1}\},
\end{equation}
increasing $C_*$ by fixed norm-convention factors if necessary.
For $\norm{J}_\infty<\varepsilon_{\rm pt}$, put
\begin{equation}\label{pc-full-perturbation}
 \Psi_J=\Theta'_J+\Pi^*\Theta_J\Pi,\qquad
 K_J=G_0\Psi_J.
\end{equation}
The actual full homogeneous flat inverse is the convergent series
\begin{equation}\label{pc-neumann-inverse}
 G_1^Jf=(I-K_J)^{-1}G_0f
                  =\sum_{n\ge0}K_J^nG_0f.
\end{equation}
It is holomorphic from the complex bounded-current ball to operators
from bounded fine forcing into $N_\beta$, uniformly for
$0\le\beta\le\beta_0$. With $g_N$ the baseline bounded-forcing-to-$N$
constant, it obeys
\begin{align}
 N_\varpi(G_1^Jf)
       &\le2g_N\norm{f}_{\infty,\varpi},
       \label{pc-inverse-N}\\
 N_{\beta,\varpi}(G_1^Jf)
       &\le(g_\beta+2C_*\varepsilon_{\rm pt}g_N)
                              \norm{f}_{\infty,\varpi}.
       \label{pc-inverse-Nbeta}
\end{align}
All constants are independent of $M=L^j$ and compatible volume.
For forcing supported in cell $m$, choosing
$\varpi(n)=e^{\kappa\operatorname{dist}(n,m)}$ gives exponential
coarse-cell locality in the pointwise target norms. The auxiliary
inverse is the same result with only $\Theta'_J$ retained.
\end{theorem}

\begin{proof}
Equations \eqref{pc-aux-N-bound} and
\eqref{pc-constraint-N-bound} give
\[
 N_{\beta,\varpi}(K_JA)
          \le C_*\norm{J}_\infty N_\varpi(A).
\]
The target dominates $N_\varpi$, so the ratio is below $1/2$
on \eqref{pc-radius}. The series converges in weighted $N$.
For every term after the zeroth, apply the stronger $N$-to-$N_\beta$
bound to its outside $K_J$ and the $N$ contraction to the remaining
factors. Summing gives \eqref{pc-inverse-N}--\eqref{pc-inverse-Nbeta}.
The coefficient $\Theta_J=T_{E_J^TJ}$ is nonlinear in $J$ but is
analytic on the auxiliary ball with the proved linear norm bound.
Uniform convergence therefore gives Banach-valued holomorphy; no
linear-current differentiation formula is presumed for the full family.

For forcing in one cell, the exponentially weighted solution is in
physical $L^2$ on the infinite lattice. It consequently equals the
independently constructed Hilbert inverse on that input. Equality
for bounded forcing follows by the absolutely summable pointwise
cell kernels and the fixed-$M$ summability of the original operator
products. The finite-volume proof is identical with periodic
distance. The minimum with $\varepsilon_H$ ensures that both
constructions use the same actual source inverse.
\end{proof}

The normalized extensions inherit the same bounds:
$E_1^J=G_1^JQ^*q_1^J$ and $E_J=G_\pi^JQ^*q_J$.
Here $Q^*$ maps bounded coarse fields to bounded fine fields, and
the established coarse Schur bounds make $q_1^J,q_J$ bounded on
weighted coarse $\ell^\infty$. Thus their output $N_\beta$ norms
are uniform on the common ball.

\subsection{What has and has not matched the source}

The input and output norms in \eqref{pc-aux-N-bound} and
\eqref{pc-constraint-N-bound} are the explicit homogeneous flat
$N,N_\beta$ norms, up to fixed equivalent conventions for localized
H\"older seminorms. The current has only a supremum bound. These
estimates and \eqref{pc-inverse-Nbeta} supply the bounded-forcing
members of the current program and their weighted versions.

No current Laplacian estimate is asserted. The source expressly
excludes that member, and the Laplacian of $DZD^*$ is an order-zero
operator whose bounded-input pointwise norm has not been supplied.
The source's separate right-divergence members with H\"older input
are not proved by this section. Such input cannot be sharply
localized by cell indicators; its further analysis needs smooth
physical cutoffs or a direct weighted global H\"older estimate.

The fixed-step corrected-average source identification remains the
one audited from the deposited transcriptions. These results are
homogeneous and flat. They do not establish nonflat source stability,
mixed retained-domain estimates, a changing nonlinear constraint
beyond its exact flat quadratic coefficient, or an interacting
Yang--Mills ultraviolet limit.

\begin{thebibliography}{9}
\bibitem{repo}
\href{https://github.com/ethanwsussman/Yang--Mills-LLM-experiment/tree/e392ef45013a581664bab76b90867170b30fe10e}
{Pinned Yang--Mills-LLM-experiment account}, commit
\texttt{e392ef45013a581664bab76b90867170b30fe10e}.
The homogeneous weak norms are Raw I:12740--12770; the bounded
current assumption is I:12951--12966; the distinct right-divergence
members are I:13601--13615; the current assertions and exclusion of
the current Laplacian member are I:13749--13804.
\bibitem{previous}
Companion research proofs: Notebook 09, \emph{The flat fine inverse
and a current neighborhood}; Notebook 11, \emph{The full flat
constraint-current family}. Source operator matching is to the pinned
deposited transcriptions. Original primary-paper arguments have not
all been independently reverified.
\end{thebibliography}
\end{document}